\section{An\'alisis y discusi\'on}

Los resultados permiten una lectura clara. Dentro de dataset, la l\'inea cl\'asica sobre EEG monocanal s\'i alcanza baselines defendibles. Fuera de dataset, el rendimiento cae de forma abrupta, en especial para apnea. Por eso la contribuci\'on principal del trabajo no es solo reportar ``qu\'e modelo gan\'o'', sino mostrar c\'omo cambian las conclusiones al pasar de \emph{accuracy} a \emph{macro-F1}, de media por fold a contraste pareado, y de validaci\'on interna a generalizaci\'on externa.

\subsection{Respuestas a las preguntas de an\'alisis}

\begin{enumerate}[leftmargin=*]
\item \textbf{¿Qu\'e combinaci\'on de preprocesamiento + \emph{features} produjo el mejor resultado y por qu\'e cree que funcion\'o mejor?} La mejor combinaci\'on para staging fue la de ventanas de 30 s, estandarizaci\'on ajustada dentro de cada fold y \feature{eeg\_*} espectrales compactas, especialmente potencia delta, potencias relativas, \emph{zero-crossing rate} y entrop\'ia espectral. En apnea funcion\'o mejor la versi\'on extendida con Hjorth, SEF50/95, entrop\'ias y curtosis. Esa mezcla funciona porque separa tanto macroestados de sue\~no como cambios finos de microestructura que suelen acompa\~nar eventos respiratorios.

\item \textbf{¿Qu\'e tan sensible fue cada modelo al desbalance de clases, especialmente en N1 o en apnea/no-apnea?} Fue muy sensible. En staging, N1 fue la clase m\'as castigada: el mejor F1 de N1 en Sleep-EDF fue apenas 0.256 con Random Forest, mientras que W super\'o 0.91 (Tabla~\ref{tab:stage_f1_class}). En apnea, el desbalance se vuelve cr\'itico en \emph{cross-dataset}: por ejemplo, MIT-BIH \(\rightarrow\) St. Vincent da una \emph{accuracy} de 0.799 con XGBoost, pero sensibilidad 0.000, un patr\'on cl\'asico de sesgo hacia la clase mayoritaria.

\item \textbf{¿Qu\'e diferencias observa entre \emph{accuracy}, \emph{macro-F1} y \emph{kappa}? ¿Alg\'un modelo parece ``bueno'' solo porque favorece clases mayoritarias?} S\'i. XGBoost fue el mejor en \emph{accuracy} y \emph{kappa} intra-dataset sobre Sleep-EDF, pero Random Forest fue mejor en \emph{macro-F1}. Eso indica que el mejor promedio global no coincide necesariamente con el mejor equilibrio entre clases. En apnea \emph{cross-dataset}, la \emph{accuracy} alta con sensibilidad cero demuestra de forma extrema que un modelo puede parecer ``bueno'' solo por predecir la clase negativa dominante.

\item \textbf{¿Qu\'e patrones aparecen en las matrices de confusi\'on? ¿Qu\'e clases se confunden con mayor frecuencia y por qu\'e?} La matriz de confusi\'on de Sleep-EDF (Figura~\ref{fig:sleep_stage_confusion}) muestra que N1 se mezcla con N2 y con vigilia. Esto tiene sentido porque N1 es una etapa transicional y, con EEG monocanal, su firma es menos estable que la de N3 o W. En apnea, la matriz de MIT-BIH confirma una cantidad importante de falsos negativos, se\~nal de que el modelo identifica mejor ausencia de evento que presencia de evento.

\item \textbf{¿El mejor modelo dentro de un dataset tambi\'en fue el mejor en \emph{cross-dataset}? Si no, ¿qu\'e indica eso sobre generalizaci\'on?} No. XGBoost domina dentro de Sleep-EDF y MIT-BIH cuando se mira rendimiento agregado interno, pero en staging \emph{cross-dataset} SVM-RBF obtuvo el mejor \emph{macro-F1} en ambas direcciones Sleep-EDF/ISRUC. Esto indica que el mejor ajuste intra-dataset no garantiza robustez fuera de dominio; probablemente algunos modelos est\'an explotando regularidades espec\'ificas del corpus fuente.

\item \textbf{¿Qu\'e tanto cambi\'o el rendimiento al pasar de \emph{epoch-wise} impl\'icito a \emph{subject-wise} estricto? ¿Existe evidencia de fuga de informaci\'on en alg\'un dise\~no preliminar?} Nuestros resultados son m\'as bajos que varios papers que reportan 0.84--0.94 de \emph{accuracy} en Sleep-EDF \cite{mousavi2019sleepeegnet,gao2023psdrf,nakamura2017complexity}. Esa brecha es consistente con el endurecimiento del protocolo: particiones \emph{subject-wise} estrictas, tuning solo en train y una disciplina uniforme entre datasets. No se presenta evidencia de fuga en la versi\'on final, pero la diferencia con la literatura sugiere que dise\~nos menos estrictos pueden inflar el rendimiento.

\item \textbf{¿Los intervalos de confianza de dos modelos se superponen? ¿Eso cambia su conclusi\'on sobre cu\'al es realmente mejor?} S\'i. En Sleep-EDF, los IC bootstrap de \emph{macro-F1} de Random Forest y XGBoost se superponen de forma importante (Tabla~\ref{tab:bootstrap_sleep_stage}), aunque el \emph{accuracy} de XGBoost queda ligeramente por encima. En MIT-BIH, los IC de AUC de Random Forest y XGBoost tambi\'en se superponen ampliamente. Esto obliga a una conclusi\'on matizada: XGBoost es el mejor compromiso global, pero no por un margen tan amplio como para descartar por completo a Random Forest.

\item \textbf{¿McNemar o la comparaci\'on de AUC apoyan una diferencia significativa o la mejora observada es solo marginal?} Para staging, McNemar respalda diferencias reales entre XGBoost y Random Forest tanto en Sleep-EDF como en MIT-BIH (Tabla~\ref{tab:stat_tests}). Para apnea en MIT-BIH tambi\'en hay diferencia significativa por McNemar (\(p=0.0043\)). Sin embargo, la diferencia bootstrap pareada de AUC favorece levemente a Random Forest sobre las predicciones agregadas. Por tanto, la mejora de XGBoost en apnea no es absoluta; depende de c\'omo se agregue la evidencia.

\item \textbf{¿Qu\'e \emph{features} resultaron m\'as importantes seg\'un RF/boosting? ¿Tienen sentido fisiol\'ogico o solo parecen estad\'isticamente \'utiles?} En staging dominaron \feature{eeg\_bandpower\_delta}, \feature{eeg\_zero\_crossing\_rate}, \feature{eeg\_rel\_power\_alpha}, \feature{eeg\_bandpower\_beta} y \feature{eeg\_spectral\_entropy}. En apnea destacaron \feature{eeg\_theta\_alpha\_ratio}, \feature{eeg\_rel\_power\_delta}, \feature{eeg\_bandpower\_beta}, \feature{eeg\_kurtosis}, \feature{eeg\_sef50} y \feature{eeg\_hjorth\_mobility}. Esas variables s\'i tienen sentido fisiol\'ogico: recogen lentificaci\'on cortical, fragmentaci\'on y cambios de estructura espectral asociados a etapas y a micro-arousals.

\item \textbf{¿C\'omo afecta el ruido, la presencia de artefactos o el cambio de base de datos al desempe\~no del modelo?} El efecto es fuerte. La ca\'ida entre intra-dataset y \emph{cross-dataset} es demasiado grande para atribuirla solo al azar. Adem\'as, no hubo un m\'odulo expl\'icito de rechazo de artefactos cl\'inicos; el pipeline conf\'ia en canalizaci\'on, remuestreo y \emph{features} robustas. Por eso el ruido residual, la diferencia de cohortes y la no equivalencia exacta del canal EEG aparecen reflejados directamente en la degradaci\'on del rendimiento.

\item \textbf{¿SVM super\'o o no a RF/boosting? Explique la raz\'on en t\'erminos de separaci\'on no lineal, tama\~no del conjunto y estabilidad.} En t\'erminos globales, no. SVM no fue el mejor baseline intra-dataset y qued\'o claramente por debajo en MIT-BIH multitarget. No obstante, en staging \emph{cross-dataset} obtuvo el mejor \emph{macro-F1}, lo que sugiere que su frontera radial sobre un conjunto compacto de \feature{eeg\_*} puede ser menos sobreajustada que la de los ensambles en ciertos dominios. Aun as\'i, RF y XGBoost fueron m\'as estables y m\'as consistentes a lo largo de la bater\'ia completa.

\item \textbf{¿Qu\'e baseline publicado es el m\'as comparable con su experimento y en qu\'e aspectos la comparaci\'on todav\'ia no es totalmente justa?} Para staging, el baseline m\'as comparable es Gao y Wu \cite{gao2023psdrf}, porque comparte Sleep-EDF, EEG monocanal y Random Forest con \emph{features} espectrales. Para apnea, el referente m\'as cercano es Barnes \emph{et al.} \cite{barnes2022apneaeegcnn}, por su restricci\'on a EEG monocanal y su foco expl\'icito en apnea. La comparaci\'on no es totalmente justa porque cambian cohortes, protocolo de partici\'on, definici\'on de labels, posibles estrategias de balance y, en algunos casos, el propio canal EEG.

\item \textbf{Si tuviera que desplegar el sistema en un escenario real con un EEG simple, ¿qu\'e modelo elegir\'ia y por qu\'e?} Elegir\'ia \textbf{XGBoost} como punto de partida, porque fue el mejor compromiso global en el paquete final: mejor \emph{accuracy}/\emph{kappa} de staging en Sleep-EDF y mejor balance medio apnea+staging en MIT-BIH. Sin embargo, no lo desplegar\'ia sin recalibraci\'on externa, ajuste de umbral y validaci\'on adicional en una cohorte destino, porque los experimentos \emph{cross-dataset} muestran que ninguno de los tres modelos es todav\'ia robusto fuera del dominio fuente.
\end{enumerate}
